\chapter{Conclusiones}\label{cap:conclusiones}

\section{Trabajo futuro}\label{cap:trabajo_futuro}

Como líneas de trabajo futuro se plantean las siguientes ampliaciones del sistema.

En primer lugar, el desarrollo de una página web para la aplicación, que permita a los usuarios acceder a determinadas funcionalidades sin necesidad de instalar la aplicación móvil.

En segundo lugar, se plantea la posibilidad de generalizar el sistema desarrollado para dar soporte no solo a la Semana Santa, sino a cualquier tipo de evento o celebración tradicional (por ejemplo, ferias, romerías o fiestas patronales). Para ello, entre otros cambios, las entidades específicas del dominio actual deberían sustituirse por clases más genéricas: las cofradías o hermandades pasarían a modelarse como \textit{organizaciones}, sus miembros o cofrades como \textit{miembros}, y las Juntas de Cofradías como \textit{conjuntos de organizaciones} encargados de coordinar a varias organizaciones dentro de un mismo evento. De esta forma, el resto del sistema (geolocalización en tiempo real, gestión de recorridos, notificaciones, etc.) podría reutilizarse prácticamente sin cambios, ya que estas funcionalidades no dependen de la naturaleza religiosa del evento, sino de su estructura organizativa. Esta generalización estaría alineada con el objetivo de escalabilidad planteado inicialmente para el proyecto (Sección~\ref{cap:objetivos}), y permitiría ampliar el alcance de la aplicación mucho más allá del ámbito de la Semana Santa.

En tercer lugar, la imagen de un \texttt{Paso} se guarda hoy como una URL de texto libre (Apéndice~\ref{cap:esquemaBD}), sin ninguna infraestructura de subida de archivos: es la Junta de Cofradías quien aloja la imagen en otro sitio y pega el enlace en el formulario, en lugar de seleccionarla directamente desde el propio dispositivo. Como mejora futura, esto podría resolverse de dos formas: guardando la imagen como un campo binario en la propia base de datos PostgreSQL --sin depender de ningún servicio externo, aunque no es la solución más eficiente para servir muchas imágenes pesadas a la vez--, o retomando el planteamiento inicial de usar \textbf{Firebase Storage} \cite{firebasestorage} --un servicio dedicado a este propósito, con mejor rendimiento a mayor escala, pero que exigiría gestionar las credenciales de un proyecto de Firebase además del backend propio ya existente--.

Por último, sería interesante adaptar la interfaz en función de la Semana Santa de cada ciudad, asociando a cada una la terminología propia que se emplea localmente, de forma que la experiencia resulte más cercana y personalizada para cada usuario.

\section{Puntos débiles}\label{cap:puntosDebiles}

El sistema presenta una seguridad limitada en el acceso de los cofrades, que se realiza mediante el escaneo de un código QR. No obstante, se ha priorizado la usabilidad frente a la seguridad, ya que el nivel de riesgo asociado a este método de acceso se considera asumible dado el contexto de uso de la aplicación.

Por otro lado, la introducción de información a través del dispositivo móvil resulta costosa para el usuario. Una alternativa más adecuada sería ofrecer una página web con capacidad de edición tanto para las juntas de cofradías como para los administradores. Sin embargo, por motivos de acotación del alcance del proyecto y de la carga de trabajo disponible, esta funcionalidad no se ha llegado a implementar.

Por otro lado, el uso de un backend propio con Spring Boot y PostgreSQL, en lugar de una plataforma \emph{Backend-as-a-Service} como Firebase, traslada a la propia autora responsabilidades que antes asumía la plataforma de forma automática: el dimensionado y escalado de la infraestructura, las copias de seguridad de la base de datos, y la aplicación de parches de seguridad al propio servidor. Railway~\cite{railway}, la plataforma elegida para el despliegue (Apéndice~\ref{cap:manual_despliegue}), simplifica buena parte de esta gestión, pero no la elimina por completo como sí ocurría con Firebase.
 
Asimismo, la posición en vivo de una procesión (RNF-08) se resuelve mediante una API REST convencional consultada periódicamente por el cliente, en lugar de mediante los \emph{listeners} en tiempo real que ofrecía Cloud Firestore de forma nativa. Esta solución cumple el requisito tal y como está definido (una actualización cada 30 segundos como mínimo), pero no ofrece la misma sensación de instantaneidad que un mecanismo de \emph{push}, y su eficiencia se degrada más rápido al aumentar el número de usuarios simultáneos que la consultan, ya que cada consulta es una petición HTTP independiente en lugar de una única notificación que el servidor reenvía a quien esté suscrito. Del mismo modo, la disponibilidad sin conexión de los datos estáticos (RNF-03), cubierta antes de forma automática por la caché local de Firestore, exige ahora una solución explícita en el cliente (Sección~\ref{cap:1_iteracion}), con el consiguiente esfuerzo de implementación adicional.
 
Por último, el desarrollo de un backend propio con una tecnología (Spring Boot) con la que la autora contaba con una experiencia previa limitada --un único proyecto de prácticas durante el Grado-- ha supuesto una curva de aprendizaje que no era necesaria con el planteamiento inicial basado en Firebase, y que se ha gestionado mediante una prueba de concepto acotada antes de comenzar la implementación completa, tal y como recoge el riesgo correspondiente (Sección~\ref{cap:analisisRiesgos}).
 
Por último, al asociarse los códigos de acceso a nivel de cofradía y no de procesión concreta (Sección~\ref{cap:1_iteracion}), un cofrade recibe notificaciones de todas las procesiones de su cofradía y no únicamente de aquellas en las que participa. Se considera una limitación menor, asumible dado que el número de procesiones por cofradía es reducido.